\section{Modality as a Condition on Assisted Analysis}
\label{sec:modality}
\draftnote{Morris; Nicole contributes Large Display and collaboration. Condense hardware exposition. Balance coverage through evidence rather than equal-length inventories.}
\subsection{Desktop/Planar Environments}
\plan{Examine precise input, linked views, dense annotation, and notebook/script integration. Ask how recommendations and delegated operations are inspected and corrected. Establish meaningful baselines for cross-environment comparisons.}
\subsection{Large Displays and Collaborative Walls}
\plan{Compare shared visibility, physical navigation, individual work regions, and group control. Examine how speech and guidance affect awareness and coordination. Distinguish individual from group assistance.}
\subsection{Virtual Reality}
\plan{Analyze spatial exploration, selection, navigation burden, and representation inspection. Identify where assistance alleviates or adds interaction cost. Tie benefit claims to tasks and studies rather than immersion alone.}
\subsection{Augmented/Mixed Reality}
\plan{Examine assistance anchored to physical context, registration, occlusion, and spatial reference. Ask how scientists detect and correct errors linking data, environment, and interpreted intent.}
\subsection{CAVE Environments}
\plan{Examine shared room-scale views, tracked viewpoints, input ownership, and collaborative control. Distinguish these conditions from head-mounted VR and non-immersive large displays.}
\subsection{Cross-Device Continuity and Comparative Synthesis}
\plan{Compare preservation of selection, representation, provenance, and control across devices. Treat hybrid workflows as combinations, not an unplanned sixth category. Identify matched comparisons and studies that cannot isolate modality from other system differences.}
